\section{Related Work}\label{sec:related-work}

\paragraph{Safety in Reasoning Models.}
Reasoning models introduce a qualitatively new failure mode: the model's own chain of thought transitions from safe deliberation to harmful generation without adversarial input~\cite{wang2025lrmsafety}. Self-jailbreak has been documented with step-level annotations~\cite{harmthoughts,selfjailbreak}; unfaithful reasoning~\cite{chen2025unfaithful,turpin2023language} motivates hidden-state monitoring. Safety-aligned models maintain refusal intentions during thinking but experience sharp drops at output generation~\cite{yin2025refusal_cliff,hagendorff2026lrm_jailbreak,hu2026refusal_dynamics,ahrend2026safer_traces}. Mitigation strategies include reasoning steering~\cite{safethink} and guardrail insertion~\cite{mao2025selfjailbreak_reasoning}; none examines how methodological choices in probing (particularly layer selection) affect safety conclusions for reasoning traces.\looseness=-1

\paragraph{Hidden-State Probing.}
Internal activations encode safety-relevant information including alignment signals~\cite{zhou2024alignment}, refusal directions~\cite{arditi2024refusal}, and separate harmfulness/refusal encodings~\cite{zhao2025harmfulness,wollschlager2025geometry}. Monitoring systems (TrajGuard~\cite{trajguard}, SafeSwitch~\cite{safeswitch}, and others~\cite{siren,safedream2026,jiang2025hiddendetect,safeint2025,rcs2026,zhou2025ghost}) assume external input manipulation at the token level; the self-jailbreak setting differs because transitions are internal and span multiple reasoning steps. Concurrently, \citet{chan2025probing} show final-layer linear probes predict alignment outcomes with strong F1, but their single-layer, single-metric evaluation does not examine whether different layers yield different conclusions. Multi-layer approaches complement our work: CLAP~\citep{clap2025} uses cross-layer attention bypassing single-layer selection, and \citet{meyoyan2026bertology} find safety features distributed across layers, echoing our observation that different metrics select different layers. These works operate at fixed layers with single metrics; we show that the choice of layer, and the metric used to select it, can reverse the conclusions drawn.\looseness=-1

\paragraph{Probing Methodology and Layer Selection.}
Linear probing spans layer-wise classifiers~\cite{alain2017understanding}, sentence probing~\cite{conneau2018probing}, and surveys~\cite{belinkov2022probing,rogers2020primer}; control tasks~\cite{hewitt2019control} distinguish structure from memorization (adopted in \cref{app:control}). Non-uniform layer-wise information distribution~\cite{tenney2019bert,marks2024geometry} and cross-layer disagreement~\cite{kaya2019shallow,zhou2020bert} are well established; our contribution is identifying a specific cascade (metric $\to$ layer $\to$ conclusion) whose severity in safety probing has not been documented. \citet{boxo2025linear_probes} show probes can rely on textual evidence; our analysis provides a complementary geometric mechanism.\looseness=-1
